\BourbakiExerciseGroup

\begin{enumerate}
\def\labelenumi{\arabic{enumi})}
\tightlist
\item
  Let \(f\) be a mapping of a topological space \(\mathbf{X}\) into a topological space \(\mathbf{X}'\) . Prove that the following statements are equivalent:
\end{enumerate}

\begin{enumerate}
\def\labelenumi{\alph{enumi})}
\item
  \(f\) is continuous on \(\mathbf{X}\) .
\item
  For every subset \(\mathbf{A}'\) of \(\mathbf{X}', f^1 (\mathring{\mathbf{A}}') \subset (f^1 (\mathbf{A}'))^0\) .
\item
  For every subset \(\mathbf{A}'\) of \(\mathbf{X}', \overline{f}^1 (\mathbf{A}') \subset \overline{f}^1 (\overline{\mathbf{A}}')\) .
\end{enumerate}

Show by an example that, if \(f\) is continuous, the sets \(\overline{f}^1 (\mathbf{A}')\) and \(\overline{f}^1 (\overline{\mathbf{A}}')\) can be distinct.

\begin{enumerate}
\def\labelenumi{\arabic{enumi})}
\setcounter{enumi}{1}
\item
  Let X, X' be two ordered sets, each topologized with the right topology (§ 1, Exercise 2). Show that a mapping \(f: \mathbf{X} \to \mathbf{X}'\) is continuous if and only if it is order-preserving.
\item
  Let X, X' be two topological spaces, f a bijection of X onto X'; then a necessary and sufficient condition for f to be a homeomorphism of X onto \(X'\) is that the topology of \(X'\) is the finest for which f is continuous.
\item
  On an ordered set X, the least upper bound of the left topology and the right topology (§ 1, Exercise 2) is the discrete topology; if X is a directed set (on the right or the left), the greatest lower bound of these two topologies is the coarsest topology on X.
\item
  On an ordered set X let \(\mathcal{C}_{0}(\mathbf{X})\) (resp. \(\mathcal{C}_{+}(\mathbf{X})\) , \(\mathcal{C}_{-}(\mathbf{X})\) ) denote the topology generated by the set of all open intervals (resp. the set of right half-open intervals, the set of left half-open intervals), bounded or not.
\end{enumerate}

\begin{enumerate}
\def\labelenumi{\alph{enumi})}
\item
  Show that if X is linearly ordered, then the open (resp. right half-open, left half-open) intervals form a base of \(\mathcal{C}_{0}(\mathbf{X})\) (resp. \(\mathcal{C}_{+}(\mathbf{X})\) , \(\mathcal{C}_{-}(\mathbf{X})\) ); and that the closed intervals are closed sets in each of these three topologies.
\item
  On a linearly ordered set X, show that the topology \(\mathcal{C}_{+}(\mathbf{X})\) is finer than the right topology (§ 1, Exercise 2), that the topology \(\mathcal{C}_{-}(\mathbf{X})\) is finer than the left topology, and that the topology \(\mathcal{C}_{0}(\mathbf{X})\) is the intersection of \(\mathcal{C}_{+}(\mathbf{X})\) and \(\mathcal{C}_{-}(\mathbf{X})\) .
\item
  Let X be a linearly ordered set. Show that a necessary and sufficient condition for \(\mathfrak{C}_{0}(\mathbf{X})\) to have a countable base is that there should be is a countable subset D of X such that, for each \(x \in X\) and each interval {]}a, b{[} containing x, there exist \(\alpha \in D\) and \(\beta \in D\) such that
\end{enumerate}

\[
a \leqslant \alpha <   x <   \beta \leqslant b
\]

{[}cf.~Chapter IV, § 2, Exercise 11 c){]}.

\begin{enumerate}
\def\labelenumi{\alph{enumi})}
\setcounter{enumi}{3}
\tightlist
\item
  Show that if X is well-ordered (Set Theory, Chapter III, § 2) then the topology \(\mathcal{G}_{+}(\mathbf{X})\) is the discrete topology, and the topologies \(\mathcal{G}_{0}(\mathbf{X})\) and \(\mathcal{G}_{-}(\mathbf{X})\) are identical.
\end{enumerate}

¶ 6) A topology on a set X is said to be quasi-maximal if it is maximal in the set of topologies in which X has no isolated points.

\begin{enumerate}
\def\labelenumi{\alph{enumi})}
\item
  Show that the following properties of a topology \(\mathcal{G}\) on X are equivalent: \(\alpha\) \(\mathcal{G}\) is quasi-maximal; \(\beta\) if X carries the topology \(\mathcal{G}\) , then X has no isolated point and every subset of X which has no isolated points is open.
\item
  Let X be a Kolmogoroff space (§ 1, Exercise 2) in which every non-empty open set is infinite. Show that there exists a quasi-maximal topology on X which is finer than the given topology (show that the set of topologies on X in which all the non-empty open sets are infinite is inductive, and use Exercise 2 d) of § 1).
\end{enumerate}

\begin{enumerate}
\def\labelenumi{\arabic{enumi})}
\setcounter{enumi}{6}
\tightlist
\item
  If X is any set, let \(\mathfrak{P}_{0}(\mathbf{X})\) denote the set of non-empty subsets of X.
\end{enumerate}

\begin{enumerate}
\def\labelenumi{\alph{enumi})}
\item
  Let X be a non-empty topological space. Let \(\mathcal{G}_{\Omega}\) (resp. \(\mathcal{G}_{\Phi}\) ) denote the coarsest topology on \(\mathfrak{P}_{0}(\mathbf{X})\) such that for every non-empty open (resp. closed) subset A of X, \(\mathfrak{P}_{0}(\mathbf{A})\) is open (resp. closed) in \(\mathfrak{P}_{0}(\mathbf{X})\) . Show that in general these two topologies are not comparable, and that the mapping \(x \to \{x\}\) is a homeomorphism of X onto a subspace of \(\mathfrak{P}_{0}(\mathbf{X})\) if this latter set carries either of the topologies \(\mathcal{G}_{\Omega}, \mathcal{G}_{\Phi}\) .
\item
  Let D be a dense subset of X. Show that the set \(\mathfrak{E}(D)\) of finite non-empty subsets of D is dense in \(\mathfrak{P}_{0}(X)\) for either of the topologies \(C_{\Omega}, C_{\Phi}\) .
\item
  If \(\mathfrak{P}_{0}(\mathbf{X})\) is ordered by inclusion, show that the topology \(\mathcal{C}_{\Omega}\) is coarser than the left topology on \(\mathfrak{P}_{0}(\mathbf{X})\) ( \(\S\) 1, Exercise 2); that \(\mathbf{A} \in \mathfrak{P}_{0}(\mathbf{X})\) is isolated in the topology \(\mathcal{C}_{\Omega}\) if and only if \(\mathbf{A} = \{x\}\) , where \(x\) is an isolated point of \(\mathbf{X}\) ; and that \(\mathbf{A} \in \mathfrak{P}_{0}(\mathbf{X})\) is isolated in the topology \(\mathcal{C}_{\Phi}\) if and only if \(\mathbf{X}\) is a finite discrete space and \(\mathbf{A} = \mathbf{X}\) .
\item
  Prove that if \(\mathfrak{P}_0(\mathbf{X})\) carries the topology \(\mathcal{C}_{\Omega}\) (resp. \(\mathcal{C}_{\Phi}\) ) and
\end{enumerate}

\[
\mathfrak {P} _ {\mathbf {0}} (\mathbf {X}) \times \mathfrak {P} _ {\mathbf {0}} (\mathbf {X})
\]

carries the topology which is the product of \(\mathcal{C}_{\Omega}\) (resp. \(\mathcal{C}_{\Phi}\) ) with itself, then the mapping \((\mathbf{M}, \mathbf{N}) \to \mathbf{M} \cup \mathbf{N}\) of \(\mathfrak{P}_0(\mathbf{X}) \times \mathfrak{P}_0(\mathbf{X})\) into \(\mathfrak{P}_0(\mathbf{X})\) is continuous.

\begin{enumerate}
\def\labelenumi{\alph{enumi})}
\setcounter{enumi}{4}
\tightlist
\item
  Show that if \(\mathfrak{P}_{0}(\mathbf{X})\) carries the topology \(\mathfrak{T}_{\Phi}\) , then the mapping \(\mathbf{M} \to \overline{\mathbf{M}}\) is continuous.
\end{enumerate}

\(f)\) Let \(\mathbf{X}\) , \(\mathbf{Y}\) be two topological spaces, \(f: \mathbf{X} \to \mathbf{Y}\) a continuous map. Show that if both \(\mathfrak{P}_{0}(\mathbf{X})\) and \(\mathfrak{P}_{0}(\mathbf{Y})\) carry the topology \(\mathcal{G}_{\Omega}\) , and if both carry the topology \(\mathcal{G}_{\Phi}\) , then the mapping \(\mathbf{M} \to f(\mathbf{M})\) of \(\mathfrak{P}_{0}(\mathbf{X})\) into \(\mathfrak{P}_{0}(\mathbf{Y})\) is continuous.

\begin{enumerate}
\def\labelenumi{\alph{enumi})}
\setcounter{enumi}{6}
\tightlist
\item
  Let X, Z be two topological spaces. Show that a mapping \(f: Z \to \mathfrak{P}_{0}(X)\) is continuous when \(\mathfrak{P}_{0}(X)\) carries the topology \(C_{\Omega}\) (resp. \(C_{\Phi}\) ) if and only if the set of all \(z \in Z\) such that \(f(z) \cap A \neq \emptyset\) is closed (resp. open) in Z for each closed (resp. open) subset A of X.
\end{enumerate}

\begin{enumerate}
\def\labelenumi{\arabic{enumi})}
\setcounter{enumi}{7}
\tightlist
\item
  Let X be a set, and c a cardinal number which is either infinite or equal to 1. The set of subsets of X consisting of X and those M ⊆ X such that card(M) \textless{} c is the set of closed subsets of X for a topology \(\mathfrak{G}_{\mathfrak{c}}\) ; \(\mathfrak{G}_{\mathfrak{1}}\) is the coarsest topology on X and \(\mathfrak{G}_{\mathfrak{c}}\) is the discrete topology if c \textgreater{} card (X). Every permutation of X is an automorphism of the topological space obtained by giving X a topology \(\mathfrak{G}_{\mathfrak{c}}\) .
\end{enumerate}

Show that conversely any topology \(\mathcal{C}\) on \(\mathbf{X}\) , which has the property that every permutation of \(\mathbf{X}\) is continuous with respect to \(\mathcal{C}\) , is necessarily one of the topologies \(\mathcal{C}_{\mathfrak{c}}\) . {[}Let \(\mathfrak{c}\) be the smallest cardinal such that card \((\mathbf{F}) < \mathfrak{c}\) for each closed subset \(\mathbf{F} \neq \mathbf{X}\) in the topology \(\mathcal{C}\) ; observe that when \(\mathbf{X}\) is infinite and \(\mathfrak{c} >\mathrm{card}(\mathbf{X})\) , \(\mathcal{C}\) is necessarily the discrete topology.{]}

\begin{enumerate}
\def\labelenumi{\arabic{enumi})}
\setcounter{enumi}{8}
\tightlist
\item
  \begin{enumerate}
  \def\labelenumii{\alph{enumii})}
  \tightlist
  \item
    With the hypotheses of Proposition 4 of no. 3, show that the coarsest topology on X for which the \(f_{t}\) are continuous is also the finest topology \(\mathfrak{G}\) on X which has the following property: every mapping \(g\) of a topological space Z into X such that the \(f_{t} \circ g\) are continuous is continuous (with respect to \(\mathfrak{G}\) ).
  \end{enumerate}
\end{enumerate}

\begin{enumerate}
\def\labelenumi{\alph{enumi})}
\setcounter{enumi}{1}
\tightlist
\item
  With the hypotheses of Proposition 6 of no. 4, show that the finest topology on X for which the \(f_{i}\) are continuous is also the coarsest topology \(\mathcal{G}\) on X which has the following property: every mapping \(g\) of X into a topological space Z such that the \(g \circ f_{i}\) are continuous is continuous (with respect to \(\mathcal{G}\) ).
\end{enumerate}

\begin{enumerate}
\def\labelenumi{\arabic{enumi}.}
\setcounter{enumi}{9}
\tightlist
\item
  On a set X, let M be a directed set of topologies without isolated points. Show that the least upper bound of M in the set of all topologies on X is a topology without isolated points.
\end{enumerate}

¶ 11. A topological space is said to be solvable if it is the disjoint union of two dense subsets.

\begin{enumerate}
\def\labelenumi{\alph{enumi})}
\item
  If a topological space X is solvable, show that X has no isolated points and that every open subspace of X is solvable. Show also conversely that if X is a topological space and B is a set of non-empty solvable open subspaces of X such that every non-empty open subset of X contains a set belonging to B, then X is solvable (consider a maximal set of mutually disjoint open sets belonging to B).
\item
  If a Kolmogoroff space X is solvable, then every non-empty open subset of X is infinite {[}use Exercise 2d) of §1{]}.
\item
  Let X be a topological space such that 1) the smallest cardinal b of a non-empty open subset of X is infinite, and 2) there is a base B of the topology of X such that card (B) ≤ b. Show that X is solvable. (Use transfinite induction to construct two disjoint dense subsets of X). By the method of Set Theory, Chapter III, § 6, Exercise 24 a). In particular, the rational line is solvable.
\item
  A topological space X which has no isolated points is said to be isodyne if card (U) = card (X) for every non-empty open subset of X. Show that in a space X with no isolated points, every non-empty open set in X contains a non-empty open subspace which is isodyne. Deduce that if every isodyne open subspace of X is soluble, then X is soluble.
\end{enumerate}

